\section{The dielectric matrix}\label{app:D.7}

This section carries out the matrix algebra behind the dielectric matrix \cref{eq:4.4-epsGG} and the macroscopic dielectric function \cref{eq:4.4-epsmac} of \cref{sec:4.4}. It covers the inversion of the matrix Dyson equation, the macroscopic average, the equivalent form of $\varepsilon_{\mathrm{mac}}$ with the long-range Coulomb term removed, the limit $\vect{q} \to 0$, and the two conventions for the matrix. The treatment follows Refs.~\cite{adler1962quantum,wiser1963dielectric} and \cite{onida2002electronic}.

\subsection*{Inversion}

Fix $\vect{q}$ in the first zone and $\omega$, and write every quantity as a matrix in the reciprocal-lattice vectors: the improper response $R = (R_{\vect{G}\vect{G}'})$ of \cref{eq:4.4-RGG}, a proper response $\tilde R$, and the diagonal Coulomb matrix $v$ with elements $v_{\vect{G}} = v_{\mathrm{C}}(|\vect{q}+\vect{G}|)$. An external potential with components $\phi_{\mathrm{ext},\vect{G}'}$ induces the density $\delta\nd_{e,\vect{G}} = \sum_{\vect{G}'}R_{\vect{G}\vect{G}'}\phi_{\mathrm{ext},\vect{G}'}$. Each component of the induced density adds the Hartree potential $v_{\vect{G}}\,\delta\nd_{e,\vect{G}}$ at its own wavevector, since Coulomb's law is diagonal in momentum, and the total potential is $\phi_{\mathrm{tot}} = \phi_{\mathrm{ext}} + v\,\delta\nd_e$ as a vector. The proper response is defined by $\delta\nd_e = \tilde R\,\phi_{\mathrm{tot}}$. Eliminating $\phi_{\mathrm{tot}}$ for arbitrary $\phi_{\mathrm{ext}}$,
\begin{equation}
  R = \tilde R + \tilde R\,v\,R, \qquad R = (\mathbf{1} - \tilde Rv)^{-1}\tilde R ,
  \label{eq:D.7-dyson}
\end{equation}
the solution of \cref{eq:C.7-dyson} with $\mathbf{X}_0 = \tilde R$ and $\mathbf{K} = v$. The dielectric matrix and its inverse are defined by $\phi_{\mathrm{ext}} = \varepsilon\,\phi_{\mathrm{tot}}$ and $\phi_{\mathrm{tot}} = \varepsilon^{-1}\phi_{\mathrm{ext}}$, so that
\begin{equation}
  \varepsilon = \mathbf{1} - v\tilde R, \qquad \varepsilon^{-1} = \mathbf{1} + vR .
  \label{eq:D.7-eps}
\end{equation}
That these are inverse matrices follows from the push-through identity $v(\mathbf{1} - \tilde Rv)^{-1} = (\mathbf{1} - v\tilde R)^{-1}v$, which holds because $(\mathbf{1} - v\tilde R)\,v = v\,(\mathbf{1} - \tilde Rv)$. With \cref{eq:D.7-dyson},
\begin{equation*}
  \mathbf{1} + vR = \mathbf{1} + (\mathbf{1} - v\tilde R)^{-1}v\tilde R = (\mathbf{1} - v\tilde R)^{-1}\bigl[(\mathbf{1} - v\tilde R) + v\tilde R\bigr] = (\mathbf{1} - v\tilde R)^{-1} .
\end{equation*}
The order of the factors matters. In element form, $\varepsilon_{\vect{G}\vect{G}'} = \delta_{\vect{G}\vect{G}'} - v_{\vect{G}}\tilde R_{\vect{G}\vect{G}'}$: the Coulomb factor carries the row index, the wavevector at which the potential is produced. This is \cref{eq:4.4-epsGG}.

\subsection*{The macroscopic average}

A field that varies slowly on the scale of the cell has, at wavevector $\vect{q} \to 0$, only the component $\vect{G}' = 0$. The total potential it produces has components at every $\vect{G}$, $\phi_{\mathrm{tot},\vect{G}} = (\varepsilon^{-1})_{\vect{G}\vect{0}}\,\phi_{\mathrm{ext},\vect{0}}$, which describe how the field inside the cell follows the electron density. The average over a cell of $\ee^{-\ii(\vect{q}+\vect{G})\cdot\vect{r}}$ vanishes for $\vect{G} \neq 0$ (\cref{eq:A.3-periodic}), so the macroscopic total potential is the component $\vect{G} = 0$ alone. The macroscopic dielectric function, the ratio of external to macroscopic total potential, is therefore
\begin{equation}
  \varepsilon_{\mathrm{mac}} = \frac{\phi_{\mathrm{ext},\vect{0}}}{\phi_{\mathrm{tot},\vect{0}}} = \frac{1}{(\varepsilon^{-1})_{\vect{0}\vect{0}}} .
  \label{eq:D.7-mac}
\end{equation}
The head $\varepsilon_{\vect{0}\vect{0}}$ answers a different question: the external potential needed at $\vect{G} = 0$ to produce a total potential that has only a $\vect{G} = 0$ component. No macroscopic external field produces such a total potential in a crystal, and $\varepsilon_{\vect{0}\vect{0}}$ is not what an experiment measures. The head of the inverse is obtained from the block inversion \cref{eq:C.7-block} with the upper-left block of size one. By \cref{eq:C.7-head},
\begin{equation}
  \varepsilon_{\mathrm{mac}} = \varepsilon_{\vect{0}\vect{0}} - \sum_{\vect{G},\vect{G}'\neq0}\varepsilon_{\vect{0}\vect{G}}\,\bigl[(\varepsilon_{\neq0})^{-1}\bigr]_{\vect{G}\vect{G}'}\,\varepsilon_{\vect{G}'\vect{0}} ,
  \label{eq:D.7-head}
\end{equation}
with $\varepsilon_{\neq0}$ the block with $\vect{G}, \vect{G}' \neq 0$. The sum is the local-field correction. It vanishes when the wings $\varepsilon_{\vect{0}\vect{G}}$ and $\varepsilon_{\vect{G}\vect{0}}$ vanish, that is, when a uniform potential induces no density modulated on the scale of the cell, as in a homogeneous system.

\subsection*{The long-range term removed}

A second form of $\varepsilon_{\mathrm{mac}}$ isolates the one element of $v$ that is singular as $\vect{q} \to 0$. Split the Coulomb matrix as $v = v_0 + \bar v$, where $v_0$ has the single nonzero element $v_{\vect{0}} = v_{\mathrm{C}}(q)$ in the head and $\bar v$ holds all the others. Define $\bar R$ by the Dyson equation with $\bar v$ alone,
\begin{equation}
  \bar R = \tilde R + \tilde R\,\bar v\,\bar R .
  \label{eq:D.7-Rbar}
\end{equation}
It is the response to the total potential with the Hartree potential at $\vect{G} = 0$ left out. Equivalently, it is the response of the electrons with the Coulomb kernel at the single wavevector $\vect{q}$ removed from their Hamiltonian. The full response follows from $\bar R$ by a second Dyson equation with $v_0$. From $R^{-1} = \tilde R^{-1} - v$ and $\bar R^{-1} = \tilde R^{-1} - \bar v$ (\cref{app:C.7}) we have $R^{-1} = \bar R^{-1} - v_0$, which is
\begin{equation*}
  R = \bar R + \bar R\,v_0\,R .
\end{equation*}
Since $v_0$ connects only the index $\vect{0}$ to itself, the head of this equation closes on itself: $(\bar Rv_0R)_{\vect{0}\vect{0}} = \bar R_{\vect{0}\vect{0}}\,v_{\mathrm{C}}(q)\,R_{\vect{0}\vect{0}}$, so $R_{\vect{0}\vect{0}} = \bar R_{\vect{0}\vect{0}}/[1 - v_{\mathrm{C}}(q)\bar R_{\vect{0}\vect{0}}]$. Then $(\varepsilon^{-1})_{\vect{0}\vect{0}} = 1 + v_{\mathrm{C}}(q)R_{\vect{0}\vect{0}} = 1/[1 - v_{\mathrm{C}}(q)\bar R_{\vect{0}\vect{0}}]$, and
\begin{equation}
  \varepsilon_{\mathrm{mac}}(\vect{q},\omega) = 1 - v_{\mathrm{C}}(q)\,\bar R_{\vect{0}\vect{0}}(\vect{q},\omega) .
  \label{eq:D.7-macbar}
\end{equation}
The macroscopic function has the form of \cref{eq:4.4-eps} for a homogeneous system, with $\bar R$ in place of $\tilde R$. The local fields are now inside $\bar R$, as the microscopic Hartree potentials $\bar v$; the exchange and correlation effects are inside $\tilde R$. Since $\bar R$ is a response of the electrons to a potential, causal and with a local Hamiltonian, the absorptive part $\varepsilon_2 = -v_{\mathrm{C}}(q)\Imag\bar R_{\vect{0}\vect{0}}$ obeys the f-sum rule \cref{eq:4.4-fsumeps} with the local fields included~\cite{onida2002electronic}. In the absorption spectrum computed from \cref{eq:D.7-macbar} the long-range term is absent, and the Coulomb singularity appears only in the loss function.

\subsection*{The limit $\vect{q} \to 0$}

At small $q$ the head $v_{\mathrm{C}}(q) = e^2/q^2$ diverges, while $v_{\vect{G}}$ stays finite for $\vect{G} \neq 0$. For an insulator the proper response is analytic at $\vect{q} = 0$, and the continuity equation fixes how its elements vanish there. Taking matrix elements of the continuity equation \cref{eq:4.3-continuity} between eigenstates, $(E_f - E_i)\langle f|n_e(\vect{q})|i\rangle = \langle f|\vect{q}\cdot\vect{j}_e(\vect{q})|i\rangle$, so $\langle f|n_e(\vect{q})|i\rangle = \vect{q}\cdot\langle f|\vect{j}_e(\vect{q})|i\rangle/(E_f - E_i)$, which is $O(q)$ because the gap keeps the denominator away from zero. Every element of $\tilde R$ with an index $\vect{0}$ carries one such factor, and
\begin{equation*}
  \tilde R_{\vect{0}\vect{0}} = O(q^2), \qquad \tilde R_{\vect{0}\vect{G}},\ \tilde R_{\vect{G}\vect{0}} = O(q), \qquad \tilde R_{\vect{G}\vect{G}'} = O(1) .
\end{equation*}
The head of the dielectric matrix, $\varepsilon_{\vect{0}\vect{0}} = 1 - (e^2/q^2)\tilde R_{\vect{0}\vect{0}}$, is then finite. Of the wings, $\varepsilon_{\vect{0}\vect{G}} = -(e^2/q^2)\tilde R_{\vect{0}\vect{G}}$ diverges as $1/q$ and $\varepsilon_{\vect{G}\vect{0}} = -v_{\vect{G}}\tilde R_{\vect{G}\vect{0}}$ vanishes as $q$, and their product in \cref{eq:D.7-head} is finite. So is $\varepsilon_{\mathrm{mac}}$, as \cref{eq:D.7-macbar} shows directly, since $\bar R$ contains no singular kernel. The finite limits depend on the direction $\hat{\vect{q}}$, through the factors $\hat{\vect{q}}\cdot\langle f|\vect{j}_e(\vect{0})|i\rangle$. This is the longitudinal projection $\hat{\vect{q}}\cdot\varepsilon\cdot\hat{\vect{q}}$ of the optical tensor of \cref{sec:4.4}, which in a cubic crystal does not depend on $\hat{\vect{q}}$ (\cref{app:C.3}). A code computes the head and wings analytically in $\hat{\vect{q}}$ from velocity matrix elements, the optical limit of \cref{app:E.6}, rather than at a small finite $q$. A metal differs: its intraband transitions have $E_f - E_i \to 0$, which makes $\tilde R_{\vect{0}\vect{0}}$ tend to a constant at $\omega = 0$, and $\varepsilon_{\mathrm{mac}}$ diverges as $1/q^2$, the perfect screening of \cref{sec:4.4}.

\subsection*{Symmetrized and unsymmetrized matrices}

Many codes and papers use the symmetrized matrix $\varepsilon^{\mathrm{s}} = \mathbf{1} - v^{1/2}\tilde R\,v^{1/2}$, with $v^{1/2}$ the diagonal matrix of $v_{\vect{G}}^{1/2}$~\cite[footnote~3]{knapen2021dark}. It is a similarity transform of the matrix used here: $\varepsilon^{\mathrm{s}} = v^{-1/2}\,\varepsilon\,v^{1/2}$, since $v^{-1/2}(v\tilde R)v^{1/2} = v^{1/2}\tilde Rv^{1/2}$. Inverting, $(\varepsilon^{\mathrm{s}})^{-1} = v^{-1/2}\varepsilon^{-1}v^{1/2}$, and its elements are $(\varepsilon^{-1})_{\vect{G}\vect{G}'}(v_{\vect{G}'}/v_{\vect{G}})^{1/2}$. The diagonal elements, which give $\varepsilon_{\mathrm{mac}}$ and the loss function at every $\vect{G}_q$, are the same in both conventions; the off-diagonal elements differ by the factor $|\vect{q}+\vect{G}|/|\vect{q}+\vect{G}'|$. The symmetrized matrix is Hermitian when $\tilde R$ is, which is convenient for inversion, and its wings are finite as $\vect{q} \to 0$. Mixing off-diagonal elements from the two conventions is an error that no diagonal check reveals.
